\documentclass{article}

\usepackage{amsmath}
\usepackage{amssymb}
\usepackage{amsthm}
\usepackage{hyperref}

\newtheorem{theorem}{Theorem}
\theoremstyle{definition}
\newtheorem{definition}{Definition}

\title{Treewidth and Colorings: a Lax paper-layer fixture}
\author{A. Author}
\date{}

\begin{document}
\maketitle

\begin{abstract}
  A small but realistic paper used to exercise the Lax marker rewriter,
  the injected \texttt{laxmark} package, and the pdf.js destination
  extractor.  It cites \cite{robertson1986graph} and
  \cite{courcelle1990monadic}.
\end{abstract}

\section{Introduction}

Graph decompositions are the standard tool for algorithmic meta-theorems;
see \cite{robertson1986graph}.  In this paper we recall the notion of
treewidth and prove a small statement about it.

We use the standard notion of
% lax begin Lax42.Colorings
proper vertex colorings
% lax end
as introduced in \cite{courcelle1990monadic}, and we will not repeat the
definition here.

\section{Treewidth}

% lax begin Lax261.Treewidth
\begin{definition}[Treewidth]
  A \emph{tree decomposition} of a graph $G = (V, E)$ is a pair
  $(T, \{X_t\}_{t \in V(T)})$ where $T$ is a tree and each $X_t \subseteq V$
  is a \emph{bag}, such that
  \begin{itemize}
    \item every vertex of $G$ lies in some bag,
    \item for every edge $uv \in E$ some bag contains both $u$ and $v$, and
    \item for every $v \in V$ the set $\{t : v \in X_t\}$ induces a subtree
      of $T$.
  \end{itemize}
  The \emph{width} of the decomposition is
  \[
    \max_{t \in V(T)} \left( |X_t| - 1 \right),
  \]
  and the \emph{treewidth} $\operatorname{tw}(G)$ is the minimum width over
  all tree decompositions of $G$.
\end{definition}
% lax end Lax261.Treewidth

The subtree condition is the one that makes dynamic programming over a
decomposition work: it guarantees that the bags containing a fixed vertex
form a connected region of the tree.

\section{A theorem with a proof}

% lax begin Lax261.Treewidth
\begin{theorem}
  \label{thm:forest}
  A graph $G$ satisfies $\operatorname{tw}(G) \le 1$ if and only if $G$ is
  a forest.
\end{theorem}

% lax begin Lax261Proofs.thm3
\begin{proof}
  If $G$ is a forest, take $T = G$ with $X_t = \{t\} \cup \{p(t)\}$ for a
  fixed rooting; every bag has size at most two, so the width is at most
  one.  Conversely, suppose $\operatorname{tw}(G) \le 1$ and let $C$ be a
  cycle in $G$.  Every bag has size at most two, and by the subtree
  condition the bags meeting $C$ form a subtree, which forces two
  non-adjacent vertices of $C$ to share a bag, a contradiction.
\end{proof}
% lax end

% lax end

\section{A marker that spans a page break}

% lax begin Lax261Proofs.longproof
The passage below is deliberately long so that its closing marker lands on
a later page than its opening marker.  The rewriter does not know or care
about pagination; only the extractor observes that the two destinations
resolve to different page indices.

Padding paragraph one.  The subtree condition is the one that makes dynamic
programming over a decomposition work, because it guarantees that the bags
containing a fixed vertex form a connected region of the tree, and therefore
a bottom-up traversal can maintain a table indexed by the intersection of a
bag with a partial solution without ever revisiting a vertex it has already
forgotten.  This is the standard argument and we repeat it only to consume
vertical space.

Padding paragraph two.  The subtree condition is the one that makes dynamic
programming over a decomposition work, because it guarantees that the bags
containing a fixed vertex form a connected region of the tree, and therefore
a bottom-up traversal can maintain a table indexed by the intersection of a
bag with a partial solution without ever revisiting a vertex it has already
forgotten.  This is the standard argument and we repeat it only to consume
vertical space.

Padding paragraph three.  The subtree condition is the one that makes dynamic
programming over a decomposition work, because it guarantees that the bags
containing a fixed vertex form a connected region of the tree, and therefore
a bottom-up traversal can maintain a table indexed by the intersection of a
bag with a partial solution without ever revisiting a vertex it has already
forgotten.  This is the standard argument and we repeat it only to consume
vertical space.

Padding paragraph four.  The subtree condition is the one that makes dynamic
programming over a decomposition work, because it guarantees that the bags
containing a fixed vertex form a connected region of the tree, and therefore
a bottom-up traversal can maintain a table indexed by the intersection of a
bag with a partial solution without ever revisiting a vertex it has already
forgotten.  This is the standard argument and we repeat it only to consume
vertical space.

Padding paragraph five.  The subtree condition is the one that makes dynamic
programming over a decomposition work, because it guarantees that the bags
containing a fixed vertex form a connected region of the tree, and therefore
a bottom-up traversal can maintain a table indexed by the intersection of a
bag with a partial solution without ever revisiting a vertex it has already
forgotten.  This is the standard argument and we repeat it only to consume
vertical space.

Padding paragraph six.  The subtree condition is the one that makes dynamic
programming over a decomposition work, because it guarantees that the bags
containing a fixed vertex form a connected region of the tree, and therefore
a bottom-up traversal can maintain a table indexed by the intersection of a
bag with a partial solution without ever revisiting a vertex it has already
forgotten.  This is the standard argument and we repeat it only to consume
vertical space.

Padding paragraph seven.  The subtree condition is the one that makes dynamic
programming over a decomposition work, because it guarantees that the bags
containing a fixed vertex form a connected region of the tree, and therefore
a bottom-up traversal can maintain a table indexed by the intersection of a
bag with a partial solution without ever revisiting a vertex it has already
forgotten.  This is the standard argument and we repeat it only to consume
vertical space.

And here the long passage ends.
% lax end

\section{A verbatim trap}

The following block must survive verbatim.  The rewriter does not detect
verbatim environments; the count check downstream is what catches this.

\begin{verbatim}
% lax begin Trap
this text is inside a verbatim block
% lax end
\end{verbatim}

An escaped percent sign is not a comment at all: 100\% of the markers in
this file are real, and the rewriter must leave that line alone because the
percent is preceded by an odd number of backslashes.

\section{A second input file}

This section lives in \texttt{section2.tex} and is pulled in with
\verb|\input|.  Its markers must get numbers that continue the numbering of
\texttt{main.tex}, because the rewriter walks main first and the remaining
\texttt{.tex} files in sorted order.

% lax begin Lax42.Colorings
\begin{definition}[Proper coloring]
  A \emph{proper $k$-coloring} of $G = (V, E)$ is a map
  $c \colon V \to \{1, \dots, k\}$ with $c(u) \neq c(v)$ for every edge
  $uv \in E$.  The chromatic number $\chi(G)$ is the least $k$ for which a
  proper $k$-coloring exists.
\end{definition}
% lax end

Combining this with Theorem~\ref{thm:forest} gives $\chi(G) \le 2$ for every
graph of treewidth at most one, as recorded in \cite{robertson1986graph}.


\bibliographystyle{plain}
\bibliography{refs}

\end{document}
